\documentclass[12pt,a4paper]{article}
\usepackage[utf8]{inputenc}
\usepackage[T1]{fontenc}
\usepackage{amsmath,amssymb,amsthm}
\usepackage{graphicx}
\usepackage{hyperref}
\usepackage{booktabs}
\usepackage{enumitem}
\usepackage[margin=1in]{geometry}

\newtheorem{theorem}{Theorem}[section]
\newtheorem{lemma}[theorem]{Lemma}
\theoremstyle{definition}
\newtheorem{definition}{Definition}[section]

\title{A Brief Note on \LaTeX{} Themes}
\author{Jane Doe}
\date{\today}

\begin{document}

\maketitle

\begin{abstract}
	This document tests a \LaTeX{} color theme by exercising various syntax
	elements common in mathematical and technical writing.
\end{abstract}

\tableofcontents
\newpage

\section{Introduction}

The \textbf{Cauchy-Schwarz inequality} states that for all vectors
$u, v$ in an inner product space:
\[
	|\langle u, v \rangle|^2 \leq \langle u, u \rangle \cdot \langle v, v \rangle.
\]

Here is an \emph{italic} and \texttt{monospace} text. We can also use
\textsuperscript{superscripts} and \textsubscript{subscripts}.

\section{Mathematical Environments}

\subsection{Inline and Display Math}

Euler's identity is $e^{i\pi} + 1 = 0$, often called the most beautiful
equation in mathematics. The quadratic formula is:
\[
	x = \frac{-b \pm \sqrt{b^2 - 4ac}}{2a}
\]

\subsection{Align Environment}

\begin{align}
	\nabla \cdot \mathbf{E}  & = \frac{\rho}{\varepsilon_0} \label{eq:gauss}                                                     \\
	\nabla \cdot \mathbf{B}  & = 0 \label{eq:no-monopole}                                                                        \\
	\nabla \times \mathbf{E} & = -\frac{\partial \mathbf{B}}{\partial t} \label{eq:faraday}                                      \\
	\nabla \times \mathbf{B} & = \mu_0 \mathbf{J} + \mu_0 \varepsilon_0 \frac{\partial \mathbf{E}}{\partial t} \label{eq:ampere}
\end{align}

These are Maxwell's equations. See Equation~\eqref{eq:gauss} for Gauss's law.

\subsection{Matrices and Cases}

A general $2 \times 2$ matrix:
\[
	A = \begin{pmatrix}
		a & b \\
		c & d
	\end{pmatrix}
\]

The determinant is $\det(A) = ad - bc$.

Piecewise function:
\[
	f(x) = \begin{cases}
		x^2 & \text{if } x \geq 0, \\
		-x  & \text{if } x < 0.
	\end{cases}
\]

\subsection{Sums, Products, Integrals}

\begin{gather}
	\sum_{k=1}^{n} k = \frac{n(n+1)}{2} \\
	\prod_{i=1}^{n} i = n! \\
	\int_0^\infty e^{-x^2} \, dx = \frac{\sqrt{\pi}}{2} \\
	\oint_{\partial \Omega} \mathbf{F} \cdot d\mathbf{r} = \iint_\Omega (\nabla \times \mathbf{F}) \cdot d\mathbf{S}
\end{gather}

\section{Theorems and Definitions}

\begin{definition}[Norm]
	A \emph{norm} on a vector space $V$ is a function $\|\cdot\|: V \to \mathbb{R}$
	satisfying:
	\begin{enumerate}[label=(\roman*)]
		\item $\|v\| \geq 0$ for all $v \in V$, with equality iff $v = 0$.
		\item $\|\alpha v\| = |\alpha| \|v\|$ for all $\alpha \in \mathbb{R}$, $v \in V$.
		\item $\|u + v\| \leq \|u\| + \|v\|$ (triangle inequality).
	\end{enumerate}
\end{definition}

\begin{theorem}[Cauchy-Schwarz]
	For all $u, v$ in an inner product space:
	\[ |\langle u, v \rangle| \leq \|u\| \cdot \|v\| \]
\end{theorem}

\begin{proof}
	If $v = 0$, both sides are zero. Otherwise, let $w = u - \frac{\langle u, v \rangle}{\langle v, v \rangle} v$.
	Then $0 \leq \|w\|^2 = \|u\|^2 - \frac{|\langle u, v \rangle|^2}{\|v\|^2}$,
	which gives the result.
\end{proof}

\section{Tables}

\begin{table}[h]
	\centering
	\caption{Comparison of methods}
	\label{tab:methods}
	\begin{tabular}{@{} l c r @{}}
		\toprule
		Method     & Accuracy & Time (ms) \\
		\midrule
		Baseline   & 85.2\%   & 120       \\
		Improved   & 91.7\%   & 95        \\
		Our method & 96.3\%   & 88        \\
		\bottomrule
	\end{tabular}
\end{table}

As shown in Table~\ref{tab:methods}, our method outperforms the alternatives.

\section{Figures and Links}

See Figure~\ref{fig:placeholder} for an illustration.

\begin{figure}[h]
	\centering
	\fbox{\parbox{0.6\textwidth}{\centering\vspace{3cm}Placeholder\vspace{3cm}}}
	\caption{A placeholder figure.}
	\label{fig:placeholder}
\end{figure}

For more details, visit \url{https://www.latex-project.org/}.

\section{Lists}

\begin{itemize}
	\item First item
	\item Second item
	      \begin{itemize}
		      \item Nested item
		      \item Another nested item
	      \end{itemize}
	\item Third item
\end{itemize}

\begin{enumerate}
	\item Step one
	\item Step two
	\item Step three
\end{enumerate}

\begin{description}
	\item[Algorithm] A finite set of unambiguous instructions.
	\item[Heuristic] A technique that finds a good enough solution.
\end{description}

\section{Conclusion}

This concludes our brief tour of \LaTeX{} features.

\end{document}
